\subsubsection{Motiviación}

Fuera del mundo de la optimización tenemos, por el momento, humanos que realizan tareas. 
Humanos que tienen sentimientos, deseos, se enfocan, se despistan, etc. Al fin y al cabo,
son los que llevan a cabo el trabajo. Es fácil pensar y parece trivial que una solución
óptima matemática tendrá una ejecución óptima en la vida real. En nuestro dominio de 
camiones recolectores de residuos de reciclaje, una persona que ya tiene una ruta definida
en la cabeza, decirle que no es una buena ruta y que debe seguir un ruta que puede resultar
laberintica y confusa, pero óptima matemáticamente hablando, puede llegar a tener efectos
adeversos en la persona quien realiza la tarea haciendo que la solución óptima no termine
siendo la óptima finalmente.

\vspace{3mm}

Por este motivo, nos proponemos perder óptimialidad sobre la ruta pero poder obtener
una mucho más interpretable y fácil de realizar. Este desafío no se traduce trivialmente
al lenguaje de modelos de programación lineal entera. Por lo cuál, exploramos vairas estrategias
para obtener el resultado desado.

\section{Distancia entre vecinos}

Dada una secuencia de aristas, que representa el recorrido que realiza un camión sobre la zona a 
recolectar. Podemos definir su belleza como la distancia que hay entre la primer pasada de una calle
con su adjacente en la secuencia. Por lo tanto, si queremos un recorrido bello vamos a querer que esta 
distnacia se lo más chica posible.

\vspace{3mm}

Vamos a plantear un modelo que minimice esta distancia. Como función objetivo vamos a querer minimizar
la sumatoria entre el módulo ...